\documentclass[11pt,a4paper]{article}
\usepackage[T1]{fontenc}
\usepackage[utf8]{inputenc}
\usepackage[polish]{babel}
\usepackage{lmodern}
\usepackage{geometry}
\usepackage{microtype}
\usepackage{booktabs}
\usepackage[hidelinks]{hyperref}
\geometry{margin=26mm}
\setlength{\parskip}{0.55em}
\setlength{\parindent}{0pt}
\title{Świadomość interoceptywna i aleksytymia jako korelaty trudności w regulacji emocji u osób dorosłych}
\author{Mateusz Pomianek}
\date{Projekt badania, wersja 1.0.0  | 28 listopada 2026}
\begin{document}
\maketitle
\begin{center}
\textit{Interoceptive Awareness and Alexithymia as Correlates of Emotion Regulation Difficulties in Adults}
\end{center}
\begin{abstract}
Projektowane przekrojowe badanie ankietowe ma oszacować związki między wymiarami samoopisowej świadomości interoceptywnej, aleksytymią i trudnościami w regulacji emocji u dorosłych. Głównym celem jest ocena związku trudności identyfikowania uczuć z trudnościami regulacji emocji po uwzględnieniu innych wymiarów aleksytymii i wybranych współzmiennych. Rozdzielone zostaną interoceptywna uwaga i interpretacja sygnałów ciała od obiektywnej trafności interoceptywnej, której ankieta nie mierzy.

\textbf{Status:} prospektywny szkic naukowy (\texttt{IDEA\_REVIEW}); nie przeprowadzono badania, nie uzyskano wyników ani zatwierdzenia etycznego.
\end{abstract}
\tableofcontents
\section{Wprowadzenie i uzasadnienie}

Garfinkel et al. (2015) rozróżnili wykonanie zadania interoceptywnego, samoopis i metapoznawczy wgląd. MAIA-2 mierzy samoopisowe aspekty interocepcji (Mehling et al., 2018). Podskala identyfikowania uczuć TAS-20 (Bagby et al., 1994) i DERS (Gratz \& Roemer, 2004) dostarczają wskaźników postulowanych konstruktów. Zbieżność treści samoopisowych ogranicza interpretację ich korelacji jako zależności między odrębnymi mechanizmami.

Interocepcja obejmuje procesy wykrywania, interpretowania i wykorzystywania sygnałów pochodzących z organizmu. Samoopisowa świadomość interoceptywna nie jest tożsama z trafnością wykonania zadań interoceptywnych. Aleksytymia jest konstruktem obejmującym między innymi trudność identyfikowania i opisywania uczuć. Regulacja emocji oznacza procesy wpływające na dynamikę przeżywania i ekspresji emocji; trudności w jej realizacji stanowią odrębny konstrukt. Ze względu na częściowe podobieństwo treści pozycji kwestionariuszowych obserwowane korelacje wymagają krytycznej oceny trafności rozróżniającej.
\section{Cel i pytania badawcze}
\paragraph{PB1.} Jaki jest związek trudności identyfikowania uczuć z deklarowanymi trudnościami regulacji emocji?

\paragraph{PB2.} Czy podskale samoopisowej świadomości interoceptywnej są niezależnie związane z trudnościami regulacji emocji?

\paragraph{PB3.} Czy interakcja pomiędzy wymiarami interocepcji i aleksytymii ma dodatkową wartość wyjaśniającą ponad efekty główne?

\section{Hipotezy i status potwierdzenia}
\paragraph{H1 (główna).} wyższy poziom trudności identyfikowania uczuć wiąże się z większymi trudnościami regulacji emocji.

\paragraph{H2 (drugorzędowa).} wybrane adaptacyjne wymiary świadomości interoceptywnej są ujemnie związane z trudnościami regulacji emocji po kontroli aleksytymii.

\paragraph{H3 (eksploracyjna).} siła związku aleksytymii z regulacją emocji zależy od wybranych wymiarów świadomości interoceptywnej.

Hipotezy są propozycjami do zatwierdzenia przed zbieraniem danych. Ich zapis nie stanowi potwierdzenia empirycznego.
\section{Projekt badania i uczestnicy}
Jednorazowe anonimowe badanie ankietowe pełnoletnich uczestników, rekrutowanych bez deklarowania reprezentatywności. Po uzyskaniu świadomej zgody: informacje demograficzne, pytania o aktualny stres i stan zdrowia w minimalnym niezbędnym zakresie, skalę interocepcji, aleksytymii i trudności regulacji emocji. Długość ankiety i kolejność narzędzi zostaną ustalone w pilotażu; można zrównoważyć kolejność skal. Należy uprzednio sprawdzić licencje, dostępność zwalidowanych polskich wersji i zasady obliczania wyników.
\subsection{Liczebność próby i precyzja}
Liczebność próby musi zostać wyznaczona przed rekrutacją na podstawie realistycznego minimalnego efektu istotnego naukowo, założonego modelu i symulacji mocy lub oczekiwanej szerokości przedziałów ufności. Żadna liczba uczestników nie jest obecnie zatwierdzona. Należy przewidzieć rezerwę na odpowiedzi niekompletne i próby wykluczone według kryteriów jakości.
\section{Narzędzia i operacjonalizacja}
Interocepcja: wymiary MAIA-2, analizowane zgodnie z potwierdzoną strukturą podskal; nie tworzyć bez uzasadnienia jednego globalnego wyniku. Aleksytymia: TAS-20 z odrębnym naciskiem na trudności identyfikowania uczuć (DIF), przy zweryfikowanych zasadach licencjonowania i interpretacji. Trudności w regulacji emocji: DERS w odpowiedniej zwalidowanej wersji i z uwzględnieniem jej struktury czynnikowej. ERQ mierzy strategie regulacji emocji i nie jest bezpośrednim zamiennikiem DERS w analizie pierwotnej. Głównym estymandem jest współczynnik regresji DERS na DIF, przy określonym zestawie zmiennych towarzyszących.
\subsection{Kontrola jakości}
Przed zbieraniem danych zdefiniować próg ukończenia, obsługę prób treningowych, wskaźniki nieuwagi i duplikatów, zasady odrzucania obserwacji oraz test wykonalności technicznej. Wykluczeń nie należy optymalizować pod pożądany wynik. Prowadzić raport liczby osób na każdym etapie, bez ujawniania danych identyfikujących.
\section{Plan analizy statystycznej}
W pierwszym kroku raportować rozkłady wyników, rzetelność (np. omega, jeśli spełnione założenia), braki danych i korelacje z przedziałami ufności. Analiza główna: regresja wieloraka wyniku DERS na podskalę DIF TAS-20 z prerejestrowanymi współzmiennymi (wiek, płeć jeśli uzasadniona, poziom odczuwanego stresu) oraz diagnostyką współliniowości. W analizach drugorzędowych wprowadzać podskale MAIA-2 indywidualnie lub w modelu wielowymiarowym, jeśli wielkość próby i współliniowość na to pozwolą. Interakcje i analizy podgrup traktować jako eksploracyjne; uwzględnić korektę na wielokrotność. Przekrojowe dane nie pozwalają na identyfikację mediacji przyczynowej.
\subsection{Raportowanie i odtwarzalność}
Raportować wielkości efektów, przedziały ufności, dokładne definicje zmiennych i liczebność analizowalną. Przewidzieć udostępnienie kodu analitycznego, słownika zmiennych i materiałów syntetycznych zamiast danych wrażliwych. Przed analizą zamrozić wersję protokołu i zależności obliczeniowe.
\section{Aspekty etyczne i ochrona danych}
Ocenić, czy projekt wymaga opinii właściwej komisji etycznej. Rekrutację rozpocząć dopiero po ustaleniu zgodności prawnej, zakresu zgody, celu i czasu retencji danych oraz procedury wycofania. Nie zapisywać danych indywidualnych w repozytorium publicznym. Nie obiecywać uczestnikom korzyści diagnostycznych.
\section{Ograniczenia i alternatywne wyjaśnienia}
Wszystkie główne zmienne będą deklaratywne, co sprzyja wspólnej wariancji metody. Nakładanie treści DIF TAS-20 i części pozycji DERS może zawyżać korelacje; należy ocenić trafność rozróżniającą. MAIA-2 nie zapewnia pomiaru obiektywnej dokładności interoceptywnej. Dobór próby i zróżnicowanie kontekstu kulturowego ograniczają uogólnianie. Uczestnikom trzeba zapewnić możliwość rezygnacji i informacje o wsparciu przy dyskomforcie emocjonalnym.
\section{Przewidywana interpretacja i dyskusja}
Zidentyfikowane związki będą wskazywać na współwystępowanie cech, nie zaś dowodzić mechanizmu kierunkowego. Względnie silna rola DIF może wynikać zarówno z procesów emocjonalnych, jak i ze wspólnej treści narzędzi. Odrębne przedstawienie poszczególnych wymiarów MAIA-2 jest ważniejsze niż interpretowanie ogólnego wyniku interocepcji.
\section{Wyniki}
Na dzień 22 listopada 2026 r. nie przedstawiono danych empirycznych. Po przeprowadzeniu badania należy zamieścić pełne wyniki, w tym wyniki niezgodne z hipotezami, dane o odrzuceniach, niepewność oszacowań oraz analizy wrażliwości.
\section{Warunki przejścia do badania}
\begin{enumerate}
\item Zweryfikować istniejącą literaturę i przyporządkować cytowania do twierdzeń.
\item Ustalić finalne narzędzia, wersje językowe, licencje i procedurę pilotażową.
\item Zatwierdzić estymand, hipotezy, liczebność próby oraz plan kontroli błędów wielokrotności.
\item Rozstrzygnąć wymogi etyczne i ochronę danych, a następnie przygotować prerejestrację.
\item Dopiero wtedy rozpocząć rekrutację i gromadzenie materiału badawczego.
\end{enumerate}
\section*{Bibliografia}

{\setlength{\parindent}{0pt}
\setlength{\parskip}{5pt}
\hangindent=1.27cm
Bagby, R. M., Parker, J. D. A., \& Taylor, G. J. (1994). The twenty-item Toronto Alexithymia Scale—I. Item selection and cross-validation of the factor structure. \textit{Journal of Psychosomatic Research, 38}(1), 23--32. \url{https://doi.org/10.1016/0022-3999(94)90005-1}\par
\hangindent=1.27cm
Garfinkel, S. N., Seth, A. K., Barrett, A. B., Suzuki, K., \& Critchley, H. D. (2015). Knowing your own heart: Distinguishing interoceptive accuracy from interoceptive awareness. \textit{Biological Psychology, 104}, 65--74. \url{https://doi.org/10.1016/j.biopsycho.2014.11.004}\par
\hangindent=1.27cm
Gratz, K. L., \& Roemer, L. (2004). Multidimensional assessment of emotion regulation and dysregulation: Development, factor structure, and initial validation of the Difficulties in Emotion Regulation Scale. \textit{Journal of Psychopathology and Behavioral Assessment, 26}(1), 41--54. \url{https://doi.org/10.1023/B:JOBA.0000007455.08539.94}\par
\hangindent=1.27cm
Mehling, W. E., Acree, M., Stewart, A., Silas, J., \& Jones, A. (2018). The Multidimensional Assessment of Interoceptive Awareness, Version 2 (MAIA-2). \textit{PLOS ONE, 13}(12), Article e0208034. \url{https://doi.org/10.1371/journal.pone.0208034}\par
}

\end{document}
